\subsection{Lie Algebra Computations}
\label{subsec:lie-algebra-computations}

In this section, we recall a well-known Lie-algebraic method for showing the approximate universality of a gate set in a purely computational manner. Here is a high level overview of the approach: Using Lie-theoretic arguments, it can be shown that given a gate set, it is approximately universal if and only if it generates the Lie algebra of the special unitary group $\SU(2^n)$. It turns out that we get the operators $-iZ_1, \ldots, -iZ_n$ for ``free'' as generators given our gadget gate definition. Now, we have three ways of generating new operators from the above: conjugation by the Clifford circuit attached to the gadget, addition, and commutation. All three of these operations can be done efficiently via the tableau representation in Definition \ref{def:tableau}, with respect to the number of qubits of the gadget. Therefore, we can compute the Lie algebra generated by the gate set, and if it is the full Lie algebra, then we have computationally verified that the gate set is approximately universal.

Note that for simplicity, we ignore global phases throughout. In other words, we work inside $\SU(2^n)$ rather than $\U(2^n)$.

\begin{definition}
\label{def:clifford-group}
The \term{$n$-qubit Clifford group} is the normalizer of the Pauli group in the unitary group:
\begin{align}
    \cC_n
    :=
    \{ C \in U(2^n) \mid C \cP_n C^\dagger = \cP_n \}.
\end{align}
Equivalently, $C \in \cC_n$ if and only if conjugation by $C$ sends Pauli operators to Pauli operators.
\end{definition}

\begin{definition}
\label{def:matrix-exponential}
For any $A \in \C^{N \times N}$, define the \term{matrix exponential} of $A$ by
\begin{align}
    e^A
    :=
    \sum_{k=0}^\infty \frac{A^k}{k!}.
\end{align}
\end{definition}

\begin{definition}
\label{def:hermitian-skew-traceless}
Let $A \in \C^{n \times n}$. We say that $A$ is \term{Hermitian} if $A^\dagger=A$, \term{skew-Hermitian} if $A^\dagger=-A$, and \term{traceless} if $\tr(A)=0$, where \term{$\tr(A)$} is the \term{trace} of $A$, defined by
\begin{align}
    \tr(A)
    =
    \sum_{j=1}^n A_{jj}.
\end{align}
\end{definition}

\begin{definition}
\label{def:su-lie-algebra}
The \term{Lie algebra of $SU(2^n)$} is
\begin{align}
    \mathfrak{su}(2^n)
    :=
    \{ A \in \C^{2^n \times 2^n}
    \mid
    A^\dagger = -A,\;
    \tr(A)=0
    \}.
\end{align}
Equivalently,
\begin{align}
    \mathfrak{su}(2^n)
    =
    \{ -iH \mid H^\dagger = H,\; \tr(H)=0 \}.
\end{align}
Thus $\mathfrak{su}(2^n)$ consists exactly of the traceless skew-Hermitian operators.
\end{definition}

\begin{fact}
\label{fact:exponential-su}
For every $U \in SU(2^n)$, there exists $A \in \mathfrak{su}(2^n)$ such that
\begin{align}
    U = e^A.
\end{align}
\end{fact}

\begin{proof}
Since $U$ is unitary, the spectral theorem gives
\begin{align}
    U
    =
    V \Diag(e^{i\theta_1},\ldots,e^{i\theta_N}) V^\dagger,
\end{align}
where $N=2^n$ and $\theta_j \in \R$ for all $j \in [N]$. Since $\det(U)=1$, we have
\begin{align}
    1
    =
    e^{i(\theta_1+\cdots+\theta_N)}.
\end{align}
Therefore $\theta_1+\cdots+\theta_N \in 2\pi \Z$. By changing one $\theta_j$ by an integer multiple of $2\pi$, we may assume
\begin{align}
    \sum_{j=1}^N \theta_j = 0.
\end{align}
Now define
\begin{align}
    A
    :=
    V \Diag(i\theta_1,\ldots,i\theta_N) V^\dagger.
\end{align}
Then $A^\dagger=-A$, $\tr(A)=i\sum_j \theta_j=0$, so $A \in \mathfrak{su}(2^n)$, and
\begin{align}
    e^A
    =
    V \Diag(e^{i\theta_1},\ldots,e^{i\theta_N}) V^\dagger
    =
    U.
\end{align}
This completes the proof.
\end{proof}

\begin{fact}
The phase-free Pauli strings
$
    P_1 \otimes \cdots \otimes P_n
$
form a basis for $\C^{2^n \times 2^n}$. In particular, every Hermitian matrix is a real linear combination of Pauli strings.
\end{fact}

\begin{fact}
\label{fact:pauli-basis-su}
An operator $A$ lies in $\mathfrak{su}(2^n)$ if and only if there exist real numbers $\alpha_P \in \R$ such that
\begin{align}
    A
    =
    -i
    \sum_{P \in \cP_n^\circ}
    \alpha_P P,
\end{align}
where
$
    \cP_n^\circ
    :=
    \{ P_1 \otimes \cdots \otimes P_n
    \mid
    P_j \in \{I,X,Y,Z\}
    \}
    \setminus \{I^{\otimes n}\}.
$
\end{fact}

\begin{proof}
First the forward direction. Since each Pauli $P$ is Hermitian, the operator
\begin{align}
    H
    :=
    \sum_{P \in \cP_n^\circ}
    \alpha_P P
\end{align}
is Hermitian. Thus $A=-iH$ is skew-Hermitian. As every non-identity Pauli has trace $0$,
\begin{align}
    \tr(A)
    =
    -i
    \sum_{P \in \cP_n^\circ}
    \alpha_P \tr(P)
    =
    0.
\end{align}
Thus $A \in \mathfrak{su}(2^n)$. Conversely, suppose $A \in \mathfrak{su}(2^n)$. Since $A$ is skew-Hermitian, $H:=iA$ is Hermitian. Every Hermitian matrix can be written as a real linear combination of Pauli strings, so
\begin{align}
    H
    =
    \beta_I I
    +
    \sum_{P \in \cP_n^\circ}
    \beta_P P
\end{align}
for some $\beta_I,\beta_P \in \R$. Taking traces gives
\begin{align}
    0
    =
    \tr(H)
    =
    \beta_I \tr(I)
    +
    \sum_{P \in \cP_n^\circ}
    \beta_P \tr(P)
    =
    2^n \beta_I.
\end{align}
Hence $\beta_I=0$ and
\begin{align}
    H
    =
    \sum_{P \in \cP_n^\circ}
    \beta_P P.
\end{align}
Since $A=-iH$, we get
\begin{align}
    A
    =
    -i
    \sum_{P \in \cP_n^\circ}
    \beta_P P
\end{align}
as desired.
\end{proof}

\begin{theorem}[Baker-Campbell-Hausdorff formula]
\label{thm:bch}
For sufficiently small $A,B \in \mathfrak{su}(2^n)$,
\begin{align}
    e^A e^B
    =
    \exp\left(
    A+B+\frac{1}{2}[A,B]
    +\frac{1}{12}[A,[A,B]]
    +\frac{1}{12}[B,[B,A]]
    +\cdots
    \right),
\end{align}
where $[A,B]=AB-BA$ is the \term{commutator} \cite{Hall2015LieGroups}.
\end{theorem}

\begin{theorem}[Trotterization]
\label{thm:trotter}
This is the Lie--Trotter formula from \cite{Trotter1959ProductSemigroups}. For $A,B \in \mathfrak{su}(2^n)$,
\begin{align}
    e^{A+B}
    =
    \lim_{k \to \infty}
    \left(
    e^{A/k} e^{B/k}
    \right)^k.
\end{align}
\end{theorem}

\begin{definition}
\label{def:available-infinitesimal-generator}
Let $\mathcal{S} \subseteq SU(2^n)$ be a gate set, and let
$
    \mathcal{G}
    :=
    \overline{\langle \mathcal{S} \rangle}
$
be the closed group generated by $\mathcal{S}$. We say that $A \in \mathfrak{su}(2^n)$ is an \term{available infinitesimal generator} if
$
    e^{tA} \in \mathcal{G}
$
for every $t \in \mathbb{R}$.
\end{definition}

\begin{fact}[Generating sums]
\label{fact:generate-sums}
If $A$ and $B$ are available infinitesimal generators, then $A+B$ is also available as an infinitesimal generator.
\end{fact}

\begin{proof}
If $e^{tA}$ and $e^{tB}$ are available for all $t \in \R$, then by Theorem~\ref{thm:trotter},
\begin{align}
    e^{t(A+B)}
    =
    \lim_{k \to \infty}
    \left(
    e^{tA/k} e^{tB/k}
    \right)^k.
\end{align}
Thus $e^{t(A+B)}$ lies in the closure of the group generated by $e^{tA}$ and $e^{tB}$.
\end{proof}

\begin{fact}[Generating commutators]
\label{fact:generate-commutators}
If $A$ and $B$ are available infinitesimal generators, then $[A,B]$ is also available as an infinitesimal generator.
\end{fact}

\begin{proof}
Using the BCH formula,
\begin{align}
    e^{\epsilon A}
    e^{\epsilon B}
    e^{-\epsilon A}
    e^{-\epsilon B}
    =
    e^{\epsilon^2[A,B]+O(\epsilon^3)}.
\end{align}
For $t >  0$, set $\epsilon=t/k^{1/2}$. Then
\begin{align}
    \left(
    e^{\epsilon A}
    e^{\epsilon B}
    e^{-\epsilon A}
    e^{-\epsilon B}
    \right)^k
    =
    \left(
    e^{(t^2/k)[A,B]+O(k^{-3/2})}
    \right)^k
    \longrightarrow
    e^{t^2[A,B]}
\end{align}
as $k \to \infty$. Since $t \in \R$ is arbitrary, this completes the proof.
\end{proof}

This gives \(e^{s[A,B]}\in\mathcal G\) for every \(s\ge0\). Repeating the same
argument with \(A\) and \(B\) interchanged gives \(e^{-s[A,B]}\in\mathcal G\).
Hence \(e^{t[A,B]}\in\mathcal G\) for every \(t\in\R\).

\begin{definition}
\label{def:lie-algebra-generated-by-gates}
Let $\mathcal{S} \subseteq SU(2^n)$ be a gate set and let $\overline{\langle \mathcal{S}\rangle}$ be the closure of the group generated by $\mathcal{S}$. The \term{Lie algebra generated by $\mathcal{S}$} is the Lie algebra of $\overline{\langle \mathcal{S}\rangle}$:
\begin{align}
    \Lie(\mathcal{S})
    :=
    \{ A \in \mathfrak{su}(2^n)
    \mid
    e^{tA} \in \overline{\langle \mathcal{S}\rangle}
    \text{ for all } t \in \R
    \}.
\end{align}
\end{definition}

\begin{theorem}
\label{thm:lie-algebra-universality}
Let $\mathcal{S} \subseteq SU(2^n)$. If
\begin{align}
    \Lie(\mathcal{S})
    =
    \mathfrak{su}(2^n),
\end{align}
then $\mathcal{S}$ is approximately universal on $n$ qubits, meaning
\begin{align}
    \overline{\langle \mathcal{S}\rangle}
    =
    SU(2^n).
\end{align}
\end{theorem}

\begin{proof}
Let \(G=\overline{\langle\mathcal S\rangle}\). By the closed subgroup theorem,
\(G\) is a matrix Lie subgroup of \(\SU(2^n)\). Its identity component \(G^\circ\)
has the same Lie algebra as \(G\). If this Lie algebra is
\(\mathfrak{su}(2^n)\), then \(G^\circ\) is the connected Lie subgroup of
\(\SU(2^n)\) with Lie algebra \(\mathfrak{su}(2^n)\). Since \(\SU(2^n)\) is
connected, \(G^\circ=\SU(2^n)\), and therefore \(G=\SU(2^n)\).
\end{proof}

\begin{definition}
\label{def:special-unitary-gadget}
Let $\Lambda$ be a $n$-wire graph gadget and $G_\Lambda$ its gadget gate. Since global phases are irrelevant, we define its special-unitary representative as follows: 
Choose any scalar \(\omega(\theta)\in\U(1)\) satisfying
\[
    \omega(\theta)^{2^n}=\det(G_\Lambda(\theta))^{-1}.
\]
Define
\[
    \widetilde G_\Lambda(\theta):=\omega(\theta)G_\Lambda(\theta)\in\SU(2^n).
\]
The choice of \(\omega(\theta)\) only changes the representative by a global
phase.
\end{definition}

\begin{definition}
\label{def:gadget-lie-algebra}
Let $\Lambda$ be a unitary gadget and $G_\Lambda$ its gadget gate. Define the closed subgroup generated by the gadget as
\begin{align}
    \mathcal{G}_\Lambda
    :=
    \overline{
    \left\langle
    \widetilde{G}_\Lambda(\theta)
    \mid
    \theta \in \R^m
    \right\rangle
    }
    \subseteq
    SU(2^n).
\end{align}
Its Lie algebra is
\begin{align}
    \mathfrak{g}_\Lambda
    :=
    \Lie(\mathcal{G}_\Lambda)
    \subseteq
    \mathfrak{su}(2^n).
\end{align}
\end{definition}

\begin{lemma}[Zero-angle Clifford]
\label{lemma:zero-angle-clifford}
Let $\Lambda$ be a unitary gadget whose measurements are all in the $XY$ plane. Then the zero-angle gadget $\widetilde{G}_\Lambda(0,\ldots,0)$ is a Clifford operation. In particular,
\begin{align}
    \widetilde{G}_\Lambda(0,\ldots,0)
    \in
    \mathcal{G}_\Lambda.
\end{align}
Multiplying by a global phase does not change membership in the
Clifford group, so \(\widetilde G_\Lambda(0,\ldots,0)\) is Clifford modulo
global phase.
\end{lemma}

\begin{proof}
At angle $0$, a $XY$-plane measurement is an $X$-basis measurement. The operations are Clifford. The final claim follows because $\widetilde{G}_\Lambda(0,\ldots,0)$ is one of the generators used in Definition~\ref{def:gadget-lie-algebra}.
\end{proof}

\begin{lemma}[Clifford inverses are available]
\label{lemma:clifford-inverses-available}
If $C \in \mathcal{G}_\Lambda$ is Clifford, then $C^\dagger \in \mathcal{G}_\Lambda$.
\end{lemma}

\begin{proof}
Modulo global phase, the Clifford group is finite, so there exists \(r\ge1\)
and a phase \(e^{i\phi}\) such that \(C^r=e^{i\phi}I\). Since
\(C\in\SU(2^n)\), this phase is a \(2^n\)-th root of unity. Hence some further
power of \(C\) is exactly \(I\). Therefore \(C^{-1}=C^{s-1}\in\mathcal G_\Lambda\)
for some \(s\ge1\).
\end{proof}

\begin{lemma}[Conjugation by available Clifford gates]
\label{lemma:conjugation-available}
Let $A \in \mathfrak{g}_\Lambda$ and let $C \in \mathcal{G}_\Lambda$ be Clifford with $C^\dagger \in \mathcal{G}_\Lambda$. Then
$
    CAC^\dagger
    \in
    \mathfrak{g}_\Lambda.
$
\end{lemma}

\begin{proof}
Since $A \in \mathfrak{g}_\Lambda$, we have $e^{tA} \in \mathcal{G}_\Lambda$ for all $t \in \R$. Since $C,C^\dagger \in \mathcal{G}_\Lambda$,
$
    C e^{tA} C^\dagger
    \in
    \mathcal{G}_\Lambda.
$
But
$
    C e^{tA} C^\dagger
    =
    e^{tCAC^\dagger}.
$
Thus $e^{tCAC^\dagger} \in \mathcal{G}_\Lambda$ for all $t \in \R$, so $CAC^\dagger \in \mathfrak{g}_\Lambda$.
\end{proof}

\begin{fact}[Allowed Lie-algebra operations]
\label{fact:allowed-lie-operations}
Inside $\mathfrak{g}_\Lambda$, we may use the following closure operations:
\begin{align}
    A,B \in \mathfrak{g}_\Lambda
    &\implies
    A+B \in \mathfrak{g}_\Lambda, \\
    A,B \in \mathfrak{g}_\Lambda
    &\implies
    [A,B] \in \mathfrak{g}_\Lambda, \\
    A\in\mathfrak g_\Lambda,\; C,C^\dagger\in\mathcal G_\Lambda
    &\implies CAC^\dagger\in\mathfrak g_\Lambda.
\end{align}
\end{fact}

The previous fact is precisely the three operations promised at the start of this subsection. Namely, with these operations and some starting operators, the hope is that our gate set will generate the full Lie algebra (Theorem \ref{thm:lie-algebra-universality}). The question remains what operators we get to start with.

\begin{proof}
The first implication is Fact~\ref{fact:generate-sums}. The second implication is Fact~\ref{fact:generate-commutators}. The third implication is Lemma~\ref{lemma:conjugation-available}.
\end{proof}

\begin{fact}[Initial $Z$-rotations]
\label{fact:initial-z-rotations}
Suppose the gadget $\Lambda$ has angle parameters allowing the implementation
\begin{align}
    \widetilde{G}_\Lambda(2t,0,\ldots,0)
    =
    C e^{-itZ_1}
\end{align}
for some Clifford $C \in \mathcal{G}_\Lambda$ and all $t \in \R$. Then
$
    -iZ_1
    \in
    \mathfrak{g}_\Lambda.
$
More generally, if for each $j \in [n]$ there is a Clifford $C_j \in \mathcal{G}_\Lambda$ such that
\begin{align}
    \widetilde G_\Lambda(2t,0,\ldots,0)=e^{-itZ_1}C
\end{align}
then
$
    -iZ_j
    \in
    \mathfrak{g}_\Lambda
$
for all $j \in [n]$.
\end{fact}

\begin{proof}
By Lemma~\ref{lemma:clifford-inverses-available}, $C^\dagger \in \mathcal{G}_\Lambda$. Therefore
\begin{align}
    C^\dagger \widetilde{G}_\Lambda(2t,0,\ldots,0)
    =
    e^{-itZ_1}
    \in
    \mathcal{G}_\Lambda
\end{align}
for all $t \in \R$. Hence $-iZ_1 \in \mathfrak{g}_\Lambda$. The same argument applies to each $j$.
\end{proof}

Note that this does not mean that we have products, e.g. $-iZ_1Z_2$ for ``free'' too, since $e^{-iZ_1}e^{-iZ_2} = e^{-i(Z_1 + Z_2)}$ is the sum rather than the product. Also, no analogous result holds for $-iX_j$ or $-iY_j$, for example, by Lemma \ref{lemma:measurement-plane-maps}.

\begin{theorem}[Gadget universality via Pauli generation]
\label{thm:gadget-universality-pauli-generation}
Let $\Lambda$ be a unitary $n$-wire gadget. If
\begin{align}
    -iP \in \mathfrak{g}_\Lambda
\end{align}
for every non-identity tensor-product Pauli $P \in \cP_n^\circ$, then $\Lambda$ is approximately universal on $n$ qubits.
\end{theorem}

\begin{proof}
By Fact~\ref{fact:pauli-basis-su}, the elements $-iP$ for non-identity Pauli strings $P$ span $\mathfrak{su}(2^n)$ as a real vector space. Since $\mathfrak{g}_\Lambda$ is closed under real linear combinations by Fact~\ref{fact:generate-sums},
$
    \mathfrak{g}_\Lambda
    =
    \mathfrak{su}(2^n).
$
Therefore $\mathcal{G}_\Lambda=SU(2^n)$ by Theorem~\ref{thm:lie-algebra-universality}, so the gate set is approximately universal.
\end{proof}

We are now ready to put everything together into an algorithm.

\begin{example}
Consider the diagonal gadget. The relevant Clifford operation is
\begin{align}
    C
    =
    \CNOT_{2,0}
    \CZ_{0,1}
    (H_0 H_1 H_2).
\end{align}
Here products are written right-to-left, so \(H_0H_1H_2\) is applied first,
then \(\CZ_{0,1}\), then \(\CNOT_{2,0}\). Starting from 
\begin{align}
    ZII,
    \qquad
    IZI,
    \qquad
    IIZ,
\end{align}
we can generate new operators using $P \mapsto CPC^\dagger$ and closure. The following pseudocode implements exactly this search.

\begin{algorithm}[H]
\caption{$3$-qubit approximate universality check for the diagonal gadget}
\label{alg:three-qubit-gadget-universality}
\begin{algorithmic}[1]
\State $C \gets \CNOT_{2,0}\CZ_{0,1}(H_0H_1H_2)$
\State $\mathrm{Queue} \gets [ZII, IZI, IIZ]$
\State $\mathrm{Found} \gets \varnothing$

\While{$\mathrm{Queue}$ is nonempty}
    \State Pop a Pauli $P$ from $\mathrm{Queue}$

    \If{$P \notin \mathrm{Found}$}
        \State Add $CPC^\dagger$ to $\mathrm{Queue}$

        \For{$R \in \mathrm{Found}$}
            \If{$P$ and $R$ anticommute}
                \State Add the phase-free Pauli underlying \(PR\) to \(\mathrm{Queue}\)
            \EndIf
        \EndFor

        \State Add $P$ to $\mathrm{Found}$
    \EndIf
\EndWhile

\If{$\mathrm{Found}$ contains all $63$ non-identity three-qubit Paulis}
    \State Conclude $\mathfrak{g}=\mathfrak{su}(8)$
    \State Conclude the gadget is approximately universal
\EndIf
\end{algorithmic}
\end{algorithm}
It turns out that the diagonal gadget fails to generate all 63 non-identity Pauli strings. On the other hand, running the above algorithm for the width-3 gadget succeeds, so the width-3 gadget is at least approximately universal. This result is stronger than it first appears, the notion of which is made precise below.

\begin{definition}
\label{def:anti-diagonal-separated}
Let \(P,Q\subseteq \mathbb Z^2\). Say that \(P\) and \(Q\) are
\term{separated} if the embedding of $P \cup Q$ into $\CGrid$ has no edges between them.
\end{definition}

\begin{lemma}
\label{lem:anti-diagonal-separation}
Let \(P\subseteq \mathbb Z^2\) be finite, and define
\begin{align}
    X
    &:=
    \max_{p,q\in P}|x(p)-x(q)|,
    &
    Y
    &:=
    \max_{p,q\in P}|y(p)-y(q)|.
\end{align}
If \(t=(A,-B)\) satisfies \(A>X\) and \(B>Y\), then for every \(m\geq 1\),
the translates \(P,P+t,\ldots,P+(m-1)t\) are pairwise separated.
\end{lemma}

\begin{proof}
Let \(i<j\), set \(d=j-i\geq 1\), and take \(p,q\in P\). The displacement
from \(p+it\) to \(q+jt\) is
\begin{align}
    q+jt-(p+it)
    =
    q-p+dt.
\end{align}
Its first coordinate is at least \(-X+dA>0\), while its second coordinate is
at most \(Y-dB<0\). Hence the two vertices are incomparable in \(\CGrid\), so
there is no edge between \(P+it\) and \(P+jt\).
\end{proof}

\begin{definition}
\label{def:k-local-computation}
A \term{\(k\)-local layer} on \(N\) qubits is a unitary
\begin{align}
    U
    =
    \bigotimes_{\alpha=1}^{s}U_\alpha,
\end{align}
where the supports of the \(U_\alpha\)'s are disjoint and each \(U_\alpha\)
acts on at most \(k\) qubits. A \term{\(k\)-local computation} is a product
of \(k\)-local layers.
\end{definition}

\begin{theorem}
\label{thm:k-universal-implies-k-local}
Let \(\Lambda\) an approximately \(k\)-qubit universal tileable graph gadget. Then every
\(k\)-local layer can be implemented inside one sufficiently large
comparability grid. Hence every finite \(k\)-local computation can be
implemented layer by layer.
\end{theorem}

\begin{proof}
Let \(U=\bigotimes_{\alpha=1}^{s}U_\alpha\) be a \(k\)-local layer. Since
\(\Lambda\) is \(k\)-qubit universal, each \(U_\alpha\) can be approximated to arbitrary accuracy by some finite tiled
computation. Choose a computation with the most finite vertices. Put all vertices of one such
computation inside a finite point set \(P\). Choose \(t=(A,-B)\) with
\begin{align}
    A>\max_{p,q\in P}|x(p)-x(q)|,
    \qquad
    B>\max_{p,q\in P}|y(p)-y(q)|.
\end{align}
By Lemma~\ref{lem:anti-diagonal-separation}, the translated regions
\(P,P+t,P+2t,\ldots\) have no edges between distinct copies. Thus the
computations for the \(U_\alpha\)'s may be placed independently along an
anti-diagonal of one large \(\CGrid\), and the resulting operation is exactly
\begin{align}
    \bigotimes_{\alpha=1}^{s}U_\alpha.
\end{align}
Repeating this construction for each layer implements any finite product of
\(k\)-local layers.
\end{proof}

Thus, we may say that the width 3 gadget is approximately universal for $3$-local layers.
\end{example}